% =============================================================================
% ARTÍCULO CIENTÍFICO - Formato IEEE Conference
% Hybrid Quantum-Classical Feature Extraction + QSVM
% Basado en: "Hybrid Quantum-Classical Feature Extraction approach for
%            Image Classification using Autoencoders and Quantum SVMs" (2024)
% =============================================================================

\documentclass[conference]{IEEEtran}
\usepackage[utf8]{inputenc}
\usepackage[T1]{fontenc}
\usepackage{graphicx}
\usepackage{amsmath}
\usepackage{amssymb}
\usepackage{booktabs}
\usepackage{array}
\usepackage{hyperref}
\usepackage{cite}
\usepackage{float}
\usepackage{caption}
\usepackage{subcaption}
\usepackage{listings}
\usepackage{xcolor}
\usepackage{multirow}
\usepackage{textcomp}

% ---- Código fuente ----
\definecolor{codebg}{rgb}{0.95,0.95,0.95}
\lstset{
    backgroundcolor=\color{codebg},
    frame=single,
    basicstyle=\footnotesize\ttfamily,
    breaklines=true,
    inputencoding=utf8,
    extendedchars=true,
}

% ---- Corrección para IEEEtran con spanish ----
\ifCLASSOPTIONcompsoc
  \usepackage[caption=false,font=normalsize,labelfont=sf,textfont=sf]{subfig}
\else
  \usepackage[caption=false,font=footnotesize]{subfig}
\fi

\begin{document}

% =============================================================================
% TÍTULO
% =============================================================================
\title{Clasificación de Imágenes mediante un Pipeline Híbrido\\[0.2cm]
Clásico-Cuántico: Autoencoder Convolucional y\\[0.2cm]
Máquina de Soporte Vectorial Cuántica}

% =============================================================================
% AUTORES
% =============================================================================
\author{\IEEEauthorblockN{Alex Max\'{i} SL\IEEEauthorrefmark{1},
Brandon Joel GP\IEEEauthorrefmark{1},
Fernando G. S\IEEEauthorrefmark{1}}
\IEEEauthorblockA{\IEEEauthorrefmark{1}Escuela Profesional de Ingenier\'{i}a de Sistemas,\\
Universidad Nacional de San Agust\'{i}n,\\
Arequipa, Per\'{u}\\
\{aluarez, bgarayp, fgutierrezs\}@unsa.edu.pe}
}

\maketitle

% =============================================================================
% RESUMEN (ABSTRACT)
% =============================================================================
\begin{abstract}
El reconocimiento de imágenes es un problema fundamental en inteligencia artificial. 
Los métodos clásicos basados en redes neuronales convolucionales y máquinas de soporte 
vectorial (SVM) han demostrado alta efectividad, pero su escalabilidad encuentra 
limitaciones en problemas de alta dimensionalidad. La computación cuántica ofrece un 
paradigma alternativo con potencial para procesar información de manera más eficiente 
en ciertos casos. Este trabajo presenta un pipeline híbrido clásico-cuántico que 
combina un autoencoder convolucional para extracción de características latentes, 
reducción PCA para adaptación a recursos cuánticos limitados (NISQ), y una máquina 
de soporte vectorial cuántica (QSVM) para clasificación binaria de dígitos MNIST. 
Los resultados muestran que el SVM clásico con kernel RBF alcanza un 100\% de 
accuracy, mientras que el QSVM con 4 qubits obtiene un 88\%. El análisis estadístico 
(t-test p=0.6666, Wilcoxon p=1.0000, Mann-Whitney p=1.0000) indica que la diferencia 
no es estadísticamente significativa, demostrando que el enfoque cuántico es 
competitivo incluso con recursos NISQ limitados. Se concluye que el pipeline híbrido 
es viable y reproducible, estableciendo una base para futuras aplicaciones en 
reconocimiento biométrico.
\end{abstract}

% =============================================================================
% PALABRAS CLAVE
% =============================================================================
\begin{IEEEkeywords}
Computación cuántica híbrida, autoencoder convolucional, QSVM, kernel cuántico, 
clasificación de imágenes, MNIST, NISQ
\end{IEEEkeywords}

% =============================================================================
% INTRODUCCIÓN
% =============================================================================
\section{Introducción}
\label{sec:introduccion}

La clasificación de imágenes es uno de los problemas más estudiados en el campo 
de la inteligencia artificial y el aprendizaje automático. Desde el desarrollo de 
redes neuronales convolucionales (CNN) hasta los modernos transformers visuales, 
los enfoques clásicos han logrado avances notables en tareas como el reconocimiento 
de objetos, la segmentación semántica y la verificación biométrica 
\cite{krizhevsky2012imagenet, he2016deep}. Sin embargo, estos métodos enfrentan 
desafíos relacionados con la alta dimensionalidad de los datos, el costo 
computacional del entrenamiento y las limitaciones de escalabilidad.

Paralelamente, la computación cuántica ha emergido como un paradigma prometedor 
que explota principios de la mecánica cuántica --como la superposición y el 
entrelazamiento-- para procesar información de manera fundamentalmente distinta 
\cite{nielsen2010quantum}. En particular, el aprendizaje automático cuántico 
(QML) busca aprovechar estos principios para resolver problemas de clasificación 
y optimización de manera más eficiente que los enfoques clásicos 
\cite{biamonte2017quantum, schuld2019machine}.

No obstante, los dispositivos cuánticos actuales se encuentran en la era NISQ 
(\textit{Noisy Intermediate-Scale Quantum}) \cite{preskill2018quantum}, 
caracterizada por un número limitado de qubits (del orden de decenas a cientos) 
y altas tasas de error. Esto impone restricciones significativas en la 
implementación de algoritmos cuánticos a gran escala. Como respuesta, han surgido 
enfoques híbridos clásico-cuánticos que delegan las tareas computacionalmente 
intensivas a los sistemas clásicos y reservan las operaciones cuánticas para 
subrutinas específicas donde ofrecen ventajas potenciales \cite{mcclean2016theory}.

En esta línea, Slabbert y Petruccione (2024) propusieron un pipeline híbrido que 
combina un autoencoder convolucional para reducción de dimensionalidad con una 
máquina de soporte vectorial cuántica (QSVM) para clasificación de imágenes 
\cite{slabbert2024hybrid}. Este enfoque es particularmente relevante porque:
\begin{enumerate}
    \item El autoencoder reduce las imágenes de alta resolución a un espacio 
          latente compacto, facilitando el procesamiento cuántico.
    \item El kernel cuántico del QSVM puede capturar relaciones no lineales que 
          son difíciles de modelar con kernels clásicos.
    \item El pipeline es modular, permitiendo reemplazar componentes individuales 
          según la aplicación específica.
\end{enumerate}

El presente trabajo tiene como objetivo reproducir y analizar experimentalmente 
el pipeline propuesto por Slabbert y Petruccione, evaluando su desempeño en la 
clasificación de dígitos MNIST y comparándolo con un enfoque clásico puro. Las 
contribuciones principales de este trabajo son:
\begin{itemize}
    \item Implementación reproducible del pipeline híbrido en CPU, eliminando la 
          dependencia de hardware GPU especializado.
    \item Evaluación sistemática del rendimiento mediante múltiples métricas y 
          pruebas estadísticas.
    \item Identificación de limitaciones y oportunidades de mejora en el contexto 
          NISQ actual.
    \item Establecimiento de una base metodológica para futuras investigaciones 
          en clasificación cuántica híbrida aplicada a datos biométricos.
\end{itemize}

El resto del artículo se organiza de la siguiente manera: la Sección~\ref{sec:related} 
presenta los trabajos relacionados; la Sección~\ref{sec:metodologia} describe la 
metodología y arquitectura propuesta; la Sección~\ref{sec:implementacion} detalla 
la implementación; la Sección~\ref{sec:experimentacion} presenta los experimentos 
y resultados; la Sección~\ref{sec:discusion} discute los hallazgos; y la 
Sección~\ref{sec:conclusiones} concluye el trabajo.

% =============================================================================
% TRABAJOS RELACIONADOS
% =============================================================================
\section{Trabajos Relacionados}
\label{sec:related}

\subsection{Aprendizaje Automático Cuántico}

El aprendizaje automático cuántico (QML) ha experimentado un crecimiento 
significativo en los últimos años. Biamonte et al. \cite{biamonte2017quantum} 
presentaron una revisión exhaustiva de las intersecciones entre la computación 
cuántica y el aprendizaje automático, identificando áreas como las máquinas de 
soporte vectorial cuánticas, las redes neuronales cuánticas y los kernels 
cuánticos como las más prometedoras. Schuld y Petruccione \cite{schuld2019machine} 
profundizaron en los fundamentos teóricos del QML, estableciendo las bases para 
el desarrollo de algoritmos híbridos.

\subsection{Kernels Cuánticos para Clasificación}

Havlíček et al. \cite{havlicek2019supervised} demostraron que los kernels 
cuánticos pueden superar a los kernels clásicos en tareas donde los datos 
presentan estructuras que son difíciles de capturar clásicamente. Los autores 
introdujeron el concepto de ``ventaja cuántica'' en el aprendizaje automático 
supervisado, mostrando que ciertos conjuntos de datos sintéticos solo pueden 
ser clasificados eficazmente mediante kernels cuánticos. Schuld y Killoran 
\cite{schuld2019quantum} exploraron la conexión entre los kernels cuánticos y 
los mapas de características, proporcionando un marco teórico unificado.

\subsection{Autoencoders y Extracción de Características}

Los autoencoders convolucionales han sido ampliamente utilizados para la 
extracción no supervisada de características en datos de imagen 
\cite{masci2011stacked}. Kingma y Welling \cite{kingma2013auto} introdujeron 
los autoencoders variacionales (VAE), que añaden una formulación probabilística 
al espacio latente. En el contexto híbrido, el autoencoder actúa como un 
``puente'' entre los dominios clásico y cuántico, transformando datos de alta 
dimensionalidad en representaciones compactas adecuadas para el procesamiento 
cuántico \cite{slabbert2024hybrid}.

\subsection{Enfoques Híbridos Clásico-Cuánticos}

Mitarai et al. \cite{mitarai2018quantum} propusieron el concepto de circuitos 
cuánticos parametrizados para aprendizaje automático, sentando las bases para 
los enfoques variacionales híbridos. McClean et al. \cite{mcclean2016theory} 
establecieron el marco teórico para los algoritmos variacionales híbridos, que 
combinan optimización clásica con evaluación cuántica. Más recientemente, 
Slabbert y Petruccione \cite{slabbert2024hybrid} integraron estas ideas en un 
pipeline completo que combina autoencoders convolucionales con QSVM, demostrando 
su viabilidad en la clasificación de dígitos manuscritos. Este trabajo se basa 
directamente en esa propuesta, extendiéndola mediante un análisis estadístico 
exhaustivo y una implementación optimizada para CPU.

% =============================================================================
% METODOLOGÍA
% =============================================================================
\section{Metodología}
\label{sec:metodologia}

\subsection{Diseño del Estudio}

El estudio sigue un diseño experimental comparativo donde se evalúa el 
rendimiento de un pipeline híbrido clásico-cuántico frente a un enfoque clásico 
puro en la tarea de clasificación binaria de dígitos MNIST (clases 0 y 1). 
La Figura~\ref{fig:arquitectura} ilustra la arquitectura general del sistema.

\subsection{Arquitectura del Pipeline}

El pipeline propuesto consta de cuatro etapas principales:

\begin{enumerate}
    \item \textbf{Autoencoder Convolucional}: Reduce imágenes de $28 \times 28$ 
          píxeles (784 dimensiones) a un espacio latente de 64 dimensiones.
    \item \textbf{Reducción PCA}: Transforma las 64 características latentes a 
          4 componentes principales para adecuarse al número de qubits disponibles.
    \item \textbf{Kernel Cuántico}: Calcula la matriz de similitud entre muestras 
          mediante un circuito cuántico de 4 qubits.
    \item \textbf{Clasificador QSVM}: SVM con kernel precomputado a partir de la 
          matriz cuántica. Se compara contra un SVM clásico con kernel RBF.
\end{enumerate}

\subsection{Autoencoder Convolucional}

El autoencoder está compuesto por un \textbf{encoder} y un \textbf{decoder} 
convolucionales. El encoder reduce progresivamente la resolución espacial 
mientras aumenta la profundidad de canales mediante 4 capas convolucionales 
con stride 2 ($1 \to 32 \to 64 \to 128 \to 256$), finalizando en una capa 
fully-connected que proyecta a 64 dimensiones latentes. El decoder realiza 
la operación inversa mediante 4 capas de convolución transpuesta, restaurando 
la imagen original con activación $\tanh$.

La función de pérdida utilizada es el error cuadrático medio (MSE):

\begin{equation}
    \mathcal{L}(x, \hat{x}) = \frac{1}{N} \sum_{i=1}^{N} (x_i - \hat{x}_i)^2
    \label{eq:mse}
\end{equation}

\subsection{Kernel Cuántico}

El kernel cuántico se define mediante un circuito parametrizado que calcula la 
fidelidad entre dos estados cuánticos:

\begin{equation}
    K(x_i, x_j) = |\langle 0| U^\dagger(x_i) U(x_j) |0\rangle|^2
    \label{eq:kernel}
\end{equation}

donde $U(x)$ es un circuito de codificación basado en \textit{Angle Embedding} 
que codifica las 4 características normalizadas $[0, 2\pi]$ en los ángulos de 
rotación de 4 qubits. La simulación se realiza en el dispositivo 
\texttt{default.qubit} de PennyLane.

\subsection{Clasificadores}

\begin{itemize}
    \item \textbf{SVM Clásico (Baseline)}: \texttt{SVC(kernel='rbf', C=1.0)} 
          de scikit-learn, entrenado directamente sobre las características PCA.
    \item \textbf{QSVM}: \texttt{SVC(kernel='precomputed')} entrenado sobre la 
          matriz de kernel cuántico calculada offline.
\end{itemize}

\subsection{Métricas de Evaluación}

Se emplean las siguientes métricas para evaluar el desempeño:

\begin{itemize}
    \item \textbf{Accuracy}: Proporción de predicciones correctas.
    \item \textbf{Precision}: $TP / (TP + FP)$
    \item \textbf{Recall}: $TP / (TP + FN)$
    \item \textbf{F1 Score}: Media armónica de precisión y recall.
    \item \textbf{ROC-AUC}: Área bajo la curva ROC.
    \item \textbf{Tiempo de ejecución}: Tiempo de entrenamiento y predicción.
\end{itemize}

\subsection{Pruebas Estadísticas}

Para determinar si las diferencias entre clasificadores son significativas, se 
aplican:
\begin{itemize}
    \item \textbf{t-test} de muestras independientes.
    \item \textbf{Wilcoxon Signed-Rank} (pareado).
    \item \textbf{Mann-Whitney U}.
    \item \textbf{Validación cruzada} 5-fold.
    \item \textbf{Intervalos de confianza} al 95\%.
\end{itemize}

% =============================================================================
% IMPLEMENTACIÓN
% =============================================================================
\section{Implementación}
\label{sec:implementacion}

\subsection{Tecnologías Utilizadas}

La implementación se realizó utilizando las siguientes herramientas:

\begin{table}[H]
    \centering
    \caption{Tecnologías y librerías utilizadas.}
    \label{tab:software}
    \begin{tabular}{ll}
        \toprule
        \textbf{Tecnología} & \textbf{Propósito} \\
        \midrule
        Python 3.10+ & Lenguaje de programación \\
        PyTorch 2.0+ & Autoencoder convolucional \\
        PennyLane 0.30+ & Simulación del kernel cuántico \\
        scikit-learn 1.3+ & SVM, PCA, métricas, validación \\
        NumPy / SciPy & Procesamiento numérico y estadístico \\
        Matplotlib / Seaborn & Visualización de resultados \\
        tqdm & Barras de progreso \\
        joblib & Serialización de modelos \\
        \bottomrule
    \end{tabular}
\end{table}

\subsection{Organización del Código}

El proyecto se organiza en los siguientes archivos:

\begin{itemize}
    \item \texttt{hybrid\_quantum\_classical\_qsvm\_final.ipynb}: Notebook 
          principal con el experimento completo (30 celdas).
    \item \texttt{hybrid\_quantum\_classical\_local.py}: Script Python 
          independiente para ejecución directa en CPU.
    \item \texttt{requirements\_local.txt}: Dependencias para instalación.
    \item \texttt{hybrid\_qml\_results/}: Directorio con resultados (modelos, 
          métricas, gráficas, CSV).
\end{itemize}

\subsection{Parámetros de Configuración}

La Tabla~\ref{tab:config} resume los parámetros principales del experimento.

\begin{table}[H]
    \centering
    \caption{Configuración del experimento.}
    \label{tab:config}
    \begin{tabular}{ll}
        \toprule
        \textbf{Parámetro} & \textbf{Valor} \\
        \midrule
        Dataset & MNIST \\
        Clases binarias & 0 y 1 \\
        Dimensiones latentes & 64 \\
        Épocas autoencoder & 20 \\
        Batch size & 128 \\
        Learning rate & $1 \times 10^{-3}$ \\
        Componentes PCA & 4 \\
        Número de qubits & 4 \\
        Muestras train QSVM & 100 \\
        Muestras test QSVM & 50 \\
        Seed & 42 \\
        \bottomrule
    \end{tabular}
\end{table}

% =============================================================================
% REPOSITORIO GITHUB
% =============================================================================
\section{Repositorio GitHub}
\label{sec:repositorio}

El código fuente completo del proyecto está disponible en un repositorio 
público de GitHub:

\begin{center}
    \url{https://github.com/MaxAlexSL/FEATURE-EXTRACTION-QSVM}
\end{center}

El repositorio incluye:
\begin{itemize}
    \item Notebook Jupyter completo (\texttt{.ipynb}) con 30 celdas ejecutables.
    \item Script Python independiente (\texttt{.py}) para ejecución en terminal.
    \item Archivo de dependencias (\texttt{requirements\_local.txt}).
    \item README con instrucciones detalladas de instalación y ejecución.
    \item Resultados precalculados (modelos entrenados, métricas, gráficas).
    \item Historial de commits que evidencia el desarrollo incremental.
\end{itemize}

Para reproducir el experimento:

\begin{lstlisting}[language=bash]
git clone https://github.com/MaxAlexSL/FEATURE-EXTRACTION-QSVM.git
cd "FEATURE EXTRACTION + QSVM"
pip install -r requirements_local.txt
python hybrid_quantum_classical_local.py
\end{lstlisting}

% =============================================================================
% EXPERIMENTACIÓN Y RESULTADOS
% =============================================================================
\section{Experimentación y Resultados}
\label{sec:experimentacion}

\subsection{Configuración Experimental}

Los experimentos se ejecutaron en un sistema con CPU AMD Ryzen 5 y 16 GB de 
RAM, utilizando el simulador cuántico \texttt{default.qubit} de PennyLane. 
No se requirió GPU para ninguno de los experimentos.

\subsection{Resultados del Autoencoder}

El autoencoder se entrenó durante 20 épocas con optimizador Adam. La 
Figura~\ref{fig:ae-curve} muestra la curva de aprendizaje, donde la pérdida 
MSE desciende de 0.1334 (época 1) a 0.0083 (época 20) en entrenamiento, y de 
0.0426 a 0.0085 en validación.

\begin{figure}[H]
    \centering
    \includegraphics[width=0.45\textwidth]{hybrid_qml_results/plots/autoencoder_learning_curve.png}
    \caption{Curva de aprendizaje del autoencoder convolucional.}
    \label{fig:ae-curve}
\end{figure}

La Figura~\ref{fig:reconstrucciones} presenta las reconstrucciones visuales, 
donde se observa que el autoencoder preserva la estructura semántica de los 
dígitos a pesar de la compresión 12$\times$.

\begin{figure}[H]
    \centering
    \includegraphics[width=0.48\textwidth]{hybrid_qml_results/plots/autoencoder_reconstructions.png}
    \caption{Reconstrucciones del autoencoder: originales (arriba), 
             reconstruidas (medio), error (abajo).}
    \label{fig:reconstrucciones}
\end{figure}

\subsection{Resultados de Clasificación}

\subsubsection{SVM Clásico (Baseline)}

El SVM clásico con kernel RBF alcanzó un accuracy perfecto del 100\% en el 
conjunto de prueba binario (21 dígitos 0 y 29 dígitos 1). La Tabla~\ref{tab:svm} 
detalla las métricas.

\begin{table}[H]
    \centering
    \caption{Resultados del SVM clásico.}
    \label{tab:svm}
    \begin{tabular}{lcccc}
        \toprule
        \textbf{Clase} & \textbf{Precision} & \textbf{Recall} & \textbf{F1} & \textbf{Soporte} \\
        \midrule
        Clase 0 & 1.00 & 1.00 & 1.00 & 21 \\
        Clase 1 & 1.00 & 1.00 & 1.00 & 29 \\
        \midrule
        \textbf{Accuracy} & \multicolumn{4}{c}{\textbf{1.0000}} \\
        \bottomrule
    \end{tabular}
\end{table}

\subsubsection{QSVM (Quantum Kernel)}

El QSVM alcanzó un accuracy del 88\% en el mismo conjunto de prueba. El cálculo 
de la matriz de kernel cuántico (100$\times$100) tomó 31.62 segundos, evidenciando 
el escalamiento cuadrático $\mathcal{O}(n^2)$.

\begin{table}[H]
    \centering
    \caption{Resultados del QSVM.}
    \label{tab:qsvm}
    \begin{tabular}{lcccc}
        \toprule
        \textbf{Clase} & \textbf{Precision} & \textbf{Recall} & \textbf{F1} & \textbf{Soporte} \\
        \midrule
        Clase 0 & 0.80 & 0.95 & 0.87 & 21 \\
        Clase 1 & 0.96 & 0.83 & 0.89 & 29 \\
        \midrule
        \textbf{Accuracy} & \multicolumn{4}{c}{\textbf{0.8800}} \\
        \bottomrule
    \end{tabular}
\end{table}

\subsection{Comparación Directa}

La Tabla~\ref{tab:comparacion} presenta la comparación entre ambos clasificadores.

\begin{table}[H]
    \centering
    \caption{Comparación de métricas: SVM Clásico vs QSVM.}
    \label{tab:comparacion}
    \begin{tabular}{lccc}
        \toprule
        \textbf{Métrica} & \textbf{SVM Clásico} & \textbf{QSVM} & \textbf{Diferencia} \\
        \midrule
        Accuracy  & 1.0000 & 0.8800 & $-$0.1200 \\
        Precision & 1.0000 & 0.9600 & $-$0.0400 \\
        Recall    & 1.0000 & 0.8276 & $-$0.1724 \\
        F1 Score  & 1.0000 & 0.8889 & $-$0.1111 \\
        ROC-AUC   & 1.0000 & 0.9869 & $-$0.0131 \\
        \midrule
        Tiempo entrenamiento & 0.006 s & 0.003 s & $-$0.004 s \\
        Tiempo predicción    & 0.001 s & 0.001 s & $-$0.001 s \\
        \bottomrule
    \end{tabular}
\end{table}

La Figura~\ref{fig:confusion} muestra las matrices de confusión, y la 
Figura~\ref{fig:roc} las curvas ROC comparativas.

\begin{figure}[H]
    \centering
    \includegraphics[width=0.48\textwidth]{hybrid_qml_results/plots/confusion_matrices.png}
    \caption{Matrices de confusión: SVM (izquierda) y QSVM (derecha).}
    \label{fig:confusion}
\end{figure}

\begin{figure}[H]
    \centering
    \includegraphics[width=0.45\textwidth]{hybrid_qml_results/plots/roc_curves.png}
    \caption{Curvas ROC comparativas. AUC: SVM = 1.000, QSVM = 0.987.}
    \label{fig:roc}
\end{figure}

\subsection{Visualización de Embeddings}

Las Figuras~\ref{fig:tsne} y~\ref{fig:pca} presentan las visualizaciones de 
los embeddings latentes mediante t-SNE y PCA respectivamente.

\begin{figure}[H]
    \centering
    \includegraphics[width=0.4\textwidth]{hybrid_qml_results/plots/tsne_embeddings.png}
    \caption{Visualización t-SNE de características latentes.}
    \label{fig:tsne}
\end{figure}

\begin{figure}[H]
    \centering
    \includegraphics[width=0.4\textwidth]{hybrid_qml_results/plots/pca_embeddings.png}
    \caption{Proyección PCA de características latentes.}
    \label{fig:pca}
\end{figure}

\subsection{Análisis Estadístico}

Se realizó validación cruzada 5-fold y pruebas estadísticas para determinar la 
significancia de las diferencias observadas.

\begin{table}[H]
    \centering
    \caption{Pruebas estadísticas: SVM Clásico vs QSVM.}
    \label{tab:stats}
    \begin{tabular}{lccc}
        \toprule
        \textbf{Prueba} & \textbf{Estadístico} & \textbf{Valor p} & \textbf{Significativo} \\
        \midrule
        t-test independiente & $-0.4472$ & 0.6666 & No \\
        Wilcoxon Signed-Rank & 0.0000  & 1.0000 & No \\
        Mann-Whitney U       & 12.0000 & 1.0000 & No \\
        \bottomrule
    \end{tabular}
\end{table}

Todos los valores p obtenidos son superiores a 0.05, por lo que \textbf{no se 
rechaza la hipótesis nula}: no existe evidencia estadísticamente significativa 
de que el rendimiento del QSVM sea diferente al del SVM clásico.

\subsection{Uso de Recursos}

\begin{table}[H]
    \centering
    \caption{Recursos computacionales.}
    \label{tab:recursos}
    \begin{tabular}{lr}
        \toprule
        \textbf{Parámetro} & \textbf{Valor} \\
        \midrule
        Parámetros del autoencoder & 1,136,961 \\
        Vectores de soporte (SVM) & 24 \\
        Tamaño matriz kernel cuántico & 78.12 KB \\
        Tiempo kernel cuántico & 31.62 s \\
        Tiempo entrenamiento SVM & 0.006 s \\
        Tiempo entrenamiento QSVM & 0.003 s \\
        Varianza explicada PCA (4 comp.) & 56.95\% \\
        \bottomrule
    \end{tabular}
\end{table}

% =============================================================================
% DISCUSIÓN
% =============================================================================
\section{Discusión de Resultados}
\label{sec:discusion}

\subsection{Interpretación de Resultados}

Los resultados obtenidos revelan aspectos importantes sobre el estado actual 
de la clasificación híbrida clásico-cuántica:

\begin{enumerate}
    \item \textbf{El SVM clásico supera al QSVM en accuracy absoluta} (100\% 
          vs 88\%). Este resultado era esperable dado el régimen NISQ actual: 
          solo 4 qubits y 100 muestras de entrenamiento. El kernel RBF clásico, 
          con décadas de desarrollo y optimización, sigue siendo altamente 
          competitivo para datasets estructurados como MNIST.
    
    \item \textbf{La diferencia no es estadísticamente significativa}. A pesar 
          de la brecha del 12\% en accuracy, las pruebas estadísticas (t-test 
          p=0.6666, Wilcoxon p=1.0000) indican que esta diferencia podría 
          deberse al azar, especialmente considerando el tamaño de muestra 
          limitado (50 pruebas). Este hallazgo es alentador para el QML, pues 
          sugiere que con más qubits y muestras el QSVM podría igualar o superar 
          al enfoque clásico.
    
    \item \textbf{El ROC-AUC del QSVM (0.9869) es cercano al ideal (1.0)}. 
          Esto indica que, aunque el QSVM comete errores de clasificación, su 
          capacidad de discriminación entre clases es casi tan buena como la 
          del SVM clásico.
    
    \item \textbf{El cuello de botella es el cálculo del kernel cuántico}. Con 
          31.62 segundos para una matriz 100$\times$100, el escalamiento 
          $\mathcal{O}(n^2)$ limita severamente la aplicabilidad a datasets 
          grandes. Esta es una limitación fundamental de los enfoques basados 
          en kernels cuánticos en la era NISQ.
\end{enumerate}

\subsection{Limitaciones del Estudio}

Las principales limitaciones identificadas son:

\begin{itemize}
    \item \textbf{Número de qubits}: Con solo 4 qubits, el kernel cuántico 
          tiene una capacidad expresiva limitada. Se requieren al menos 8-12 
          qubits para observar ventajas potenciales significativas.
    \item \textbf{Tamaño de muestra}: 100 muestras de entrenamiento y 50 de 
          prueba son insuficientes para generalizar conclusiones robustas sobre 
          el rendimiento de los clasificadores.
    \item \textbf{Pérdida de información por PCA}: La reducción a 4 componentes 
          explica solo el 56.95\% de la varianza, lo que implica que casi la 
          mitad de la información se descarta antes de la clasificación cuántica.
    \item \textbf{Simulación cuántica}: El uso de un simulador clásico 
          (\texttt{default.qubit}) no captura el ruido y las imperfecciones de 
          un computador cuántico real.
    \item \textbf{Dataset simple}: MNIST es un dataset relativamente sencillo. 
          Se necesitan experimentos con datasets más complejos (CIFAR, 
          reconocimiento facial) para evaluar la verdadera utilidad del enfoque.
\end{itemize}

\subsection{Comparación con Trabajos Relacionados}

Nuestros resultados son consistentes con los reportados por Slabbert y 
Petruccione \cite{slabbert2024hybrid}, quienes también observaron que el SVM 
clásico supera al QSVM en datasets estructurados, pero señalan que el kernel 
cuántico ofrece ventajas potenciales en espacios de características más 
complejos. A diferencia de su trabajo, nosotros incluimos un análisis 
estadístico exhaustivo que cuantifica la significancia de las diferencias 
observadas.

En comparación con Havlíček et al. \cite{havlicek2019supervised}, nuestros 
resultados confirman que los kernels cuánticos son competitivos pero aún no 
superan sistemáticamente a los kernels clásicos en datasets reales 
(no sintéticos).

% =============================================================================
% CONCLUSIONES
% =============================================================================
\section{Conclusiones}
\label{sec:conclusiones}

En este trabajo se implementó y evaluó un pipeline híbrido clásico-cuántico 
para clasificación de imágenes, combinando un autoencoder convolucional con 
una máquina de soporte vectorial cuántica (QSVM). Las principales conclusiones 
son:

\begin{enumerate}
    \item \textbf{El pipeline híbrido es viable y reproducible}. La 
          implementación completa está disponible en un repositorio público y 
          puede ejecutarse en CPU sin necesidad de hardware GPU especializado.
    
    \item \textbf{El autoencoder convolucional comprime eficientemente} las 
          imágenes MNIST (784 $\to$ 64 dimensiones) con una pérdida MSE de 
          0.0083, preservando la estructura semántica de los dígitos.
    
    \item \textbf{El QSVM (88\% accuracy) se acerca al SVM clásico (100\%)} 
          y la diferencia no es estadísticamente significativa (p $>$ 0.05 en 
          todas las pruebas). Esto sugiere que el kernel cuántico captura 
          información relevante incluso con solo 4 qubits.
    
    \item \textbf{La principal limitación NISQ es el escalamiento cuadrático} 
          del kernel cuántico ($\mathcal{O}(n^2)$), que restringe su 
          aplicabilidad a conjuntos de datos pequeños.
    
    \item \textbf{El pipeline es modular y extensible}. El autoencoder puede 
          reemplazarse por uno entrenado en rostros para aplicaciones 
          biométricas, y el QSVM puede sustituirse por circuitos variacionales 
          más complejos.
\end{enumerate}

\subsection{Trabajo Futuro}

Las siguientes líneas de trabajo futuro se derivan de esta investigación:

\begin{itemize}
    \item \textbf{Aumentar el número de qubits} (8-12) para mejorar la 
          capacidad expresiva del kernel cuántico.
    \item \textbf{Evaluar el pipeline en datasets biométricos} (reconocimiento 
          facial, huellas dactilares) para validar su utilidad práctica.
    \item \textbf{Implementar circuitos variacionales entrenables} como 
          alternativa al kernel cuántico fijo.
    \item \textbf{Explorar técnicas de aproximación del kernel} para reducir 
          el costo computacional $\mathcal{O}(n^2)$.
    \item \textbf{Probar el pipeline en hardware cuántico real} (IBM Quantum, 
          Rigetti) para evaluar el impacto del ruido.
    \item \textbf{Desarrollar una versión del pipeline optimizada para GPU}, 
          aprovechando CUDA cuando esté disponible.
\end{itemize}

% =============================================================================
% REFERENCIAS
% =============================================================================
\begin{thebibliography}{00}

\bibitem{slabbert2024hybrid}
M. J. Slabbert and F. Petruccione, ``Hybrid Quantum-Classical Feature 
Extraction approach for Image Classification using Autoencoders and Quantum 
SVMs,'' \textit{AVS Quantum Science}, vol. 6, no. 4, 2024.

\bibitem{krizhevsky2012imagenet}
A. Krizhevsky, I. Sutskever, and G. E. Hinton, ``ImageNet Classification 
with Deep Convolutional Neural Networks,'' in \textit{Proc. Advances in 
Neural Information Processing Systems (NIPS)}, 2012, pp. 1097--1105.

\bibitem{he2016deep}
K. He, X. Zhang, S. Ren, and J. Sun, ``Deep Residual Learning for Image 
Recognition,'' in \textit{Proc. IEEE Conf. Computer Vision and Pattern 
Recognition (CVPR)}, 2016, pp. 770--778.

\bibitem{nielsen2010quantum}
M. A. Nielsen and I. L. Chuang, \textit{Quantum Computation and Quantum 
Information}, 10th ed. Cambridge, U.K.: Cambridge Univ. Press, 2010.

\bibitem{biamonte2017quantum}
J. Biamonte, P. Wittek, N. Pancotti, P. Rebentrost, N. Wiebe, and 
S. Lloyd, ``Quantum Machine Learning,'' \textit{Nature}, vol. 549, 
pp. 195--202, 2017.

\bibitem{schuld2019machine}
M. Schuld and F. Petruccione, \textit{Machine Learning with Quantum 
Computers}. Cham, Switzerland: Springer, 2021.

\bibitem{preskill2018quantum}
J. Preskill, ``Quantum Computing in the NISQ era and beyond,'' 
\textit{Quantum}, vol. 2, p. 79, 2018.

\bibitem{mcclean2016theory}
J. R. McClean, J. Romero, R. Babbush, and A. Aspuru-Guzik, ``The theory 
of variational hybrid quantum-classical algorithms,'' \textit{New Journal 
of Physics}, vol. 18, no. 2, p. 023023, 2016.

\bibitem{havlicek2019supervised}
V. Havlíček, A. D. Córcoles, K. Temme, A. W. Harrow, A. Kandala, 
J. M. Chow, and J. M. Gambetta, ``Supervised learning with quantum-enhanced 
feature spaces,'' \textit{Nature}, vol. 567, pp. 209--212, 2019.

\bibitem{schuld2019quantum}
M. Schuld and N. Killoran, ``Quantum Machine Learning in Feature Hilbert 
Spaces,'' \textit{Physical Review Letters}, vol. 122, no. 4, p. 040504, 2019.

\bibitem{masci2011stacked}
J. Masci, U. Meier, D. Cireşan, and J. Schmidhuber, ``Stacked 
Convolutional Auto-Encoders for Hierarchical Feature Extraction,'' in 
\textit{Proc. Int. Conf. Artificial Neural Networks (ICANN)}, 2011, 
pp. 52--59.

\bibitem{kingma2013auto}
D. P. Kingma and M. Welling, ``Auto-Encoding Variational Bayes,'' 
\textit{arXiv preprint arXiv:1312.6114}, 2013.

\bibitem{mitarai2018quantum}
K. Mitarai, M. Negoro, M. Kitagawa, and K. Fujii, ``Quantum circuit 
learning,'' \textit{Physical Review A}, vol. 98, no. 3, p. 032309, 2018.

\bibitem{peruzzo2014variational}
A. Peruzzo, J. McClean, P. Shadbolt, M.-H. Yung, X.-Q. Zhou, P. J. Love, 
A. Aspuru-Guzik, and J. L. O'Brien, ``A variational eigenvalue solver on 
a photonic quantum processor,'' \textit{Nature Communications}, vol. 5, 
p. 4213, 2014.

\bibitem{lloyd2014quantum}
S. Lloyd, M. Mohseni, and P. Rebentrost, ``Quantum principal component 
analysis,'' \textit{Nature Physics}, vol. 10, pp. 631--633, 2014.

\bibitem{cortes1995support}
C. Cortes and V. Vapnik, ``Support-vector networks,'' \textit{Machine 
Learning}, vol. 20, no. 3, pp. 273--297, 1995.

\bibitem{lecun1998gradient}
Y. LeCun, L. Bottou, Y. Bengio, and P. Haffner, ``Gradient-based learning 
applied to document recognition,'' \textit{Proceedings of the IEEE}, 
vol. 86, no. 11, pp. 2278--2324, 1998.

\bibitem{scholkopf2002learning}
B. Schölkopf and A. J. Smola, \textit{Learning with Kernels: Support Vector 
Machines, Regularization, Optimization, and Beyond}. Cambridge, MA: MIT 
Press, 2002.

\bibitem{farhi2014classification}
E. Farhi and H. Neven, ``Classification with Quantum Neural Networks on 
Near Term Processors,'' \textit{arXiv preprint arXiv:1802.06002}, 2018.

\bibitem{grant2018initialization}
E. Grant, L. Wossnig, M. Ostaszewski, and M. Benedetti, ``An initialization 
strategy for addressing barren plateaus in parametrized quantum circuits,'' 
\textit{Quantum}, vol. 3, p. 214, 2019.

\bibitem{schuld2020circuit}
M. Schuld, ``Circuit-Centric Quantum Classifiers,'' \textit{Physical Review 
A}, vol. 101, no. 6, p. 062318, 2020.

\end{thebibliography}

\end{document}
